\noindent\textit{Sửa từ đề cũ số 3. Biên soạn: Trang.}

\begin{bt}[2,0 điểm]
Thực hiện phép tính (bằng cách hợp lí nếu có thể):
\begin{listEX}[2]
\item $\dfrac{1}{2}-\dfrac{2}{3}:\dfrac{5}{6}$
\item $15\cdot\left(\dfrac{2}{3}\right)^{2}+\left(\dfrac{23}{6}\right)^{0}\cdot\dfrac{24}{16}-2\cdot\dfrac{2}{3}$
\item $\dfrac{3}{4}\cdot\dfrac{2}{9}-\dfrac{3}{4}\cdot\dfrac{11}{9}+0{,}75$
\item $\dfrac{5^{23}\cdot 2^{78}}{125^{8}\cdot 9^{13}}$
\end{listEX}
\loigiai{
a) $-\dfrac{3}{10}$.\quad
b) $\dfrac{41}{6}$.\quad
c) $0$.\quad
d) $\dfrac{2^{78}}{5\cdot 3^{26}}$.
}
\end{bt}

\begin{bt}[2,0 điểm]
Tìm số hữu tỉ $x$, biết:
\begin{listEX}[2]
\item $-2-x=(-0{,}4)^{2}$
\item $\left(\dfrac{1}{3}-x\right)\left(-\dfrac{1}{2}\right)^{5}=\left(-\dfrac{1}{2}\right)^{7}$
\item $-\dfrac{2}{3}:x=-\dfrac{7}{12}-\dfrac{5}{8}$
\item $\left|3x-\dfrac{14}{5}\right|+24\cdot\sqrt{\dfrac{25}{64}}=2^{4}$
\end{listEX}
\loigiai{
a) $x=-\dfrac{54}{25}$.\quad
b) $x=\dfrac{1}{12}$.\quad
c) $x=\dfrac{16}{29}$ (vì $x\neq 0$).\quad
d) $x=\dfrac{3}{5}$ hoặc $x=\dfrac{19}{15}$.
}
\end{bt}

\begin{bt}[1,0 điểm]
Hai người đi mua gạo. Người thứ nhất mua gạo nếp, người thứ hai mua gạo tẻ. Giá gạo tẻ rẻ hơn giá gạo nếp $20\%$. Biết khối lượng gạo tẻ người thứ hai mua nhiều hơn khối lượng gạo nếp $20\%$. Hỏi người nào trả ít tiền hơn, ít hơn bao nhiêu phần trăm so với người kia?
\loigiai{
Người mua gạo tẻ trả ít hơn $4\%$.\\
Gọi giá và khối lượng gạo nếp lần lượt là $a$ và $m$. Tiền gạo nếp là $am$, tiền gạo tẻ là $0{,}8a\cdot 1{,}2m=0{,}96am$.
}
\end{bt}

\begin{bt}[1,0 điểm]
Người ta cần hút hết nước trong một cái bể chứa đầy nước để dọn sạch bể. Lần thứ nhất hút $\dfrac{1}{3}$ lượng nước trong bể. Lần thứ hai hút $\dfrac{2}{3}$ lượng nước còn lại trong bể. Lần thứ ba hút đúng $96\,\mathrm{m}^{3}$ nước thì cạn bể. Hỏi khi bể cạn, ta cần bao nhiêu mét khối nước để làm bể đầy?
\loigiai{
$432\,\mathrm{m}^{3}$.\\
Sau lần một còn $\dfrac{2}{3}$ bể. Sau lần hai còn $\dfrac{1}{3}\cdot\dfrac{2}{3}=\dfrac{2}{9}$ bể. Gọi thể tích bể là $x$ ($\mathrm{m}^{3}$), $x>0$, ta có $\dfrac{2}{9}x=96$.
}
\end{bt}

\begin{bt}[1,5 điểm]
Cho hình vẽ bên. Biết $\widehat{A_{1}}=40^{\circ}$, $\widehat{B_{1}}=50^{\circ}$, $\widehat{AHB}=90^{\circ}$. Chứng minh $m\parallel n$.

\begin{center}
\begin{tikzpicture}[scale=1.15,line join=round,line cap=round]
	\coordinate (A) at (2.4,2.2);
	\coordinate (H) at ($(A)+1.75*(-0.7660,-0.6428)$);
	\coordinate (B) at ($(H)+1.65*(0.6428,-0.7660)$);
	\draw ($(A)+(-2.5,0)$) -- ($(A)+(1.5,0)$) node[right] {$m$};
	\draw ($(B)+(-2.2,0)$) -- ($(B)+(1.8,0)$) node[right] {$n$};
	\coordinate (ML) at ($(A)+(-1.1,0)$);
	\coordinate (NL) at ($(B)+(-1.1,0)$);
	\draw (A) -- (H) -- (B);
	\node[above] at (A) {$A$};
	\node[left] at (H) {$H$};
	\node[below] at (B) {$B$};
	\pic[draw,angle radius=5mm,angle eccentricity=1.55,"$40^{\circ}$"] {angle = ML--A--H};
	\pic[draw,angle radius=5mm,angle eccentricity=1.6,"$50^{\circ}$"] {angle = H--B--NL};
	\pic[draw,angle radius=3mm] {right angle = A--H--B};
\end{tikzpicture}
\end{center}
\loigiai{
$m\parallel n$.\\
Qua $H$ kẻ đường thẳng $d\parallel m$. Góc giữa $d$ và $AH$ bằng $40^{\circ}$. Vì $\widehat{AHB}=90^{\circ}$ nên góc giữa $d$ và $HB$ bằng $50^{\circ}$, bằng $\widehat{B_{1}}$. Suy ra $d\parallel n$. Mà $d\parallel m$ nên $m\parallel n$.
}
\end{bt}

\begin{bt}[2,0 điểm]
Cho tam giác $ABC$ cân tại $A$ và $\widehat{BAC}=80^{\circ}$. Gọi $M$ là trung điểm của $BC$. Trên tia đối của tia $MA$ lấy điểm $D$ sao cho $MD=MA$.
\begin{listEX}[1]
\item Tính $\widehat{ABC}$ và $\widehat{ACB}$.
\item Chứng minh $\triangle ABM=\triangle ACM$ và $AM\perp BC$.
\item Chứng minh $BC$ là đường trung trực của $AD$. Suy ra $\triangle ABD$ cân tại $B$ và tính $\widehat{ADB}$.
\end{listEX}
\loigiai{
a) $\widehat{ABC}=\widehat{ACB}=50^{\circ}$.\\
b) $AB=AC$, $BM=CM$, $AM$ chung nên $\triangle ABM=\triangle ACM$ (c.c.c). Hai góc tại $M$ bằng nhau và kề bù nên mỗi góc bằng $90^{\circ}$. Vậy $AM\perp BC$.\\
c) $M$ là trung điểm của $AD$ và $BC\perp AD$ nên $BC$ là đường trung trực của $AD$, do đó $DB=AB$. Tam giác $ABD$ cân tại $B$ và $\widehat{ADB}=40^{\circ}$.\\
\textit{Hướng dẫn ý c.} Chứng minh $\triangle ABM=\triangle DBM$ (c.g.c) để có $\widehat{BDM}=\widehat{BAM}=40^{\circ}$.
}
\end{bt}

\begin{bt}[0,5 điểm]
Tính giá trị của biểu thức
\[
S=\dfrac{1}{1\cdot 3}+\dfrac{1}{3\cdot 5}+\dfrac{1}{5\cdot 7}+\cdots+\dfrac{1}{23\cdot 25}.
\]
\loigiai{
$S=\dfrac{12}{25}$.\\
\textit{Hướng dẫn.} $\dfrac{1}{n(n+2)}=\dfrac{1}{2}\left(\dfrac{1}{n}-\dfrac{1}{n+2}\right)$. Tổng các số hạng rút gọn còn $\dfrac{1}{2}\left(1-\dfrac{1}{25}\right)$.
}
\end{bt}
